%------------------------------------------------------------------------------%

\usepackage{calculo_varias_variaveis-1}

\title {Fórmula de Euler}

%------------------------------------------------------------------------------%
\begin{document}

\addtitle

%------------------------------------------------------------------------------%
\section{Cálculo com Números Complexos}

%------------------------------------------------------------------------------%
\begin{frame}{Derivadas e Integrais}
  Quando derivamos ou integramos por uma variável real, $x\in\R$
  \pause
  \begin{itemize}
    \item A constante \emph{\(i\)} se comporta como uma \emph{constante}
  \end{itemize}~\\

  \pause
  Para derivar ou integrar por uma variável complexa
  \pause
  \begin{itemize}
    \item Precisamos de uma nova teoria
    \pause
    --\emph{ Variáveis Complexas}
  \end{itemize}
\end{frame}


%------------------------------------------------------------------------------%
\section{Funções de Números Complexos}

%------------------------------------------------------------------------------%
\begin{frame}{Funções de Variáveis Complexas}
  Conhecemos as Séries de Taylor das funções em $\R$
  \pause
  \[
    \sin x = x - \frac{x^3}{3!} + \frac{x^5}{5!} + \frac{x^7}{7!} - \cdots
  \]
  \pause
  \[
    \cos x = 1 - \frac{x^2}{2!} + \frac{x^4}{4!} + \frac{x^8}{8!} - \cdots
  \]
  \pause
  \[
    e^x = 1 + x + \frac{x^2}{2!} + \frac{x^3}{3!} + \frac{x^4}{4!} +  \cdots
  \]
\end{frame}

%------------------------------------------------------------------------------%
\begin{frame}{Funções de Variáveis Complexas}
  Podemos estender a definições das funções para \(\mathbb{C}\)
  \pause
  \[
    \sin z = z - \frac{z^3}{3!} + \frac{z^5}{5!} + \frac{z^7}{7!} - \cdots
  \]
  \pause
  \[
    \cos z = 1 - \frac{z^2}{2!} + \frac{z^4}{4!} + \frac{z^8}{8!} - \cdots
  \]
  \pause
  \[
    e^z = 1 + z + \frac{z^2}{2!} + \frac{z^3}{3!} + \frac{z^4}{4!} +  \cdots
  \]
\end{frame}

%------------------------------------------------------------------------------%
\section{Exponencial Complexa}

%------------------------------------------------------------------------------%
\begin{proofframe}{Justificativa}
  Escolhendo $z=i\theta$
  \begin{align*}
    \action<+->{e^{i\theta}}
    \action<+->{& = 1 + (i\theta)
      + \frac{(i\theta)^2}{2!}
      + \frac{(i\theta)^3}{3!}
      + \frac{(i\theta)^4}{4!}
      + \frac{(i\theta)^5}{5!}
      + \frac{(i\theta)^6}{6!}
      + \cdots \\[2mm]}
    \action<+->{& = 1 +  i\theta
      + \frac{i^2\theta^2}{2!}
      + \frac{i^3\theta^3}{3!}
      + \frac{i^4\theta^4}{4!}
      + \frac{i^5\theta^5}{5!}
      + \frac{i^6\theta^6}{6!}
      + \cdots \\[2mm]}
    \action<+->{& = 1 +  i\theta
      - \frac{\theta^2}{2!}
      - i\frac{\theta^3}{3!}
      + \frac{\theta^4}{4!}
      + i\frac{\theta^5}{5!}
      - \frac{\theta^6}{6!}
      + \cdots \\[2mm]}
    \action<+->{& = \left( 1 - \frac{\theta^2}{2!}
      + \frac{\theta^4}{4!}
      - \frac{\theta^6}{6!}
      + \cdots
      \right)
      + i \left( \theta - \frac{\theta^3}{3!}
      + \frac{\theta^5}{5!}
      - \cdots
      \right) \\[2mm]}
    \action<+->{& = \cos\theta + i\sin\theta}
  \end{align*}
\end{proofframe}

%------------------------------------------------------------------------------%
\begin{frame}{Identidade de Euler}
  \[
    e^{i\theta} = \cos\theta + i\sin\theta
  \]
  \pause
  ``Todas as Constantes da Matemática''
  \[
    e^{i\pi}
    \pause
    \;=\; \cos\pi + i\sin\pi
    \pause
    \;=\; -1 + i0
    \pause
    \;=\; -1
  \]
  \[
    \pause
    e^{i\pi} + 1 = 0
  \]
\end{frame}

%------------------------------------------------------------------------------%
\begin{frame}{Exponencial Complexa}
  \[
    e^{z}
    \;=\; e^{a+bi}                                   \pause
    \;=\; e^a\; e^{bi}                               \pause
    \;=\; e^a \; \big[ \cos(b) + i\sin(b)\; \big]
  \]
  \pause
  \vspace{-\baselineskip}
  \begin{itemize}[<+->]
    \item \(e^a\)                \tabto{35mm} crescimento ou decrescimento
    \item \(\cos(b) + i\sin(b)\) \tabto{35mm} oscilação
  \end{itemize}
\end{frame}

%------------------------------------------------------------------------------%
\section{Lista Mínima}

%------------------------------------------------------------------------------%
\begin{frame}{Lista Mínima}
%%  Cálculo Vol. 2 do Thomas 12\textsuperscript{a} ed. -- Seção 14.8
%%  \begin{enumerate}
%%    \item Estudar o texto da seção
%%    \item Resolver os exercícios: 3, 5, 7, 12, 16, 18, 24 e 30
%%  \end{enumerate}

  \vfill
  Atenção: A prova é baseada no livro, não nas apresentações
\end{frame}

%------------------------------------------------------------------------------%
\end{document}
%------------------------------------------------------------------------------%
